%====================================================
% PART C
% COBIT 2019
%====================================================

\chapter{COBIT 2019: IT Governance and Management Framework}

This chapter examines COBIT 2019, developed by ISACA, as a framework for the governance and management of enterprise IT. It explores the six core principles underpinning the framework, the critical distinction between governance and management, a deep dive into two selected management domains, and the performance management system based on capability levels.

%====================================================
\section{Core Principles}
%====================================================

\subsection{Principle 1: Provide Stakeholder Value}

% Meaning in the context of IT governance:

% How it supports business objectives:

% Practical example:

%---------------------------------------

\subsection{Principle 2: Holistic Approach}

% Meaning in the context of IT governance:

% How it supports business objectives:

% Practical example:

%---------------------------------------

\subsection{Principle 3: Dynamic Governance System}

% Meaning in the context of IT governance:

% How it supports business objectives:

% Practical example:

%---------------------------------------

\subsection{Principle 4: Governance Distinct from Management}

% Meaning in the context of IT governance:

% How it supports business objectives:

% Practical example:

%---------------------------------------

\subsection{Principle 5: Tailored to Enterprise Needs}

% Meaning in the context of IT governance:

% How it supports business objectives:

% Practical example:

%---------------------------------------

\subsection{Principle 6: End-to-End Governance System}

% Meaning in the context of IT governance:

% How it supports business objectives:

% Practical example:

%---------------------------------------

\subsection{COBIT 2019 versus COBIT 5}

% Explain what changed and what improvements were introduced in the 2019 revision

%%%%%%%%%%%%%%%%%%%%%%%%%%%%%%%%%%%%%%%%%%%%%%%%%%%%%

\section{Governance and Management}

%====================================================

\subsection{Governance versus Management in COBIT 2019}

% Governance: EDM domain (Evaluate, Direct, Monitor)
% Management: APO, BAI, DSS, MEA domains

%---------------------------------------

\subsection{Stakeholder Responsibilities}

\subsubsection{Governance Stakeholders}

% Board of directors, executive leadership

\subsubsection{Management Stakeholders}

% CIO, IT management, department heads

%---------------------------------------

\subsection{EDM Governance Objectives}

\subsubsection{EDM01: Ensured Governance Framework Setting and Maintenance}

\subsubsection{EDM02: Ensured Benefits Delivery}

\subsubsection{EDM03: Ensured Risk Optimisation}

\subsubsection{EDM04: Ensured Resource Optimisation}

\subsubsection{EDM05: Ensured Stakeholder Engagement}

%---------------------------------------

\subsection{Importance of the Governance and Management Separation}

% Why this separation is critical for organizational accountability

%%%%%%%%%%%%%%%%%%%%%%%%%%%%%%%%%%%%%%%%%%%%%%%%%%%%%

\section{Management Domains Deep Dive}

%====================================================

% Select TWO domains from: APO, BAI, DSS, MEA
% Replace [Domain Name] placeholders with your chosen domains

\subsection{Domain One: [Domain Name]}

\subsubsection{Purpose and Scope}

\subsubsection{Management Objectives}

% List all management objectives within this domain

\subsubsection{Selected Objective One}

% Detail and implementation example

\subsubsection{Selected Objective Two}

% Detail and implementation example

\subsubsection{Selected Objective Three}

% Detail and implementation example

\subsubsection{Contribution to IT Governance and Risk Management}

\subsubsection{Implementation Challenges}

%---------------------------------------

\subsection{Domain Two: [Domain Name]}

\subsubsection{Purpose and Scope}

\subsubsection{Management Objectives}

% List all management objectives within this domain

\subsubsection{Selected Objective One}

% Detail and implementation example

\subsubsection{Selected Objective Two}

% Detail and implementation example

\subsubsection{Selected Objective Three}

% Detail and implementation example

\subsubsection{Contribution to IT Governance and Risk Management}

\subsubsection{Implementation Challenges}

%%%%%%%%%%%%%%%%%%%%%%%%%%%%%%%%%%%%%%%%%%%%%%%%%%%%%

\section{Performance Management and Maturity}

%====================================================

\subsection{Capability Levels}

\subsubsection{Level 0: Incomplete}

\subsubsection{Level 1: Performed}

\subsubsection{Level 2: Managed}

\subsubsection{Level 3: Established}

\subsubsection{Level 4: Predictable}

\subsubsection{Level 5: Optimising}

%---------------------------------------

\subsection{Government Ministry Improvement Plan}

% Scenario: Ministry assessed at Capability Level 2
% Target: Level 4 within 18 months using the DSS domain as example

% Improvement Plan Table goes here (see templates/roadmap.tex)

%---------------------------------------

\subsection{Governance Performance Dashboard}

% Minimum 8 Key Performance Indicators (KPIs)

% KPI Table goes here (see templates/kpi-dashboard.tex)

% Minimum 6 Key Goal Indicators (KGIs)

% KGI Table goes here (see templates/kpi-dashboard.tex)

%---------------------------------------

\subsection{Alignment with CMMI and ISO 15504}

% How COBIT 2019's performance management system aligns with CMMI and ISO 15504

%%%%%%%%%%%%%%%%%%%%%%%%%%%%%%%%%%%%%%%%%%%%%%%%%%%%%

\section*{Chapter Summary}

Summarize the key findings from this chapter and briefly transition to the comparative analysis and integrated case study in the next chapter.
